\documentclass[10pt,a4paper]{article}
\usepackage[T2A]{fontenc}
\usepackage[utf8]{inputenc}
\usepackage[russian]{babel}
\usepackage[margin=1.9cm]{geometry}
\usepackage{booktabs,array,enumitem,hyperref}
\setlist{nosep,leftmargin=1.2em}
\setlength{\parindent}{0pt}\setlength{\parskip}{3pt}
\title{Поиск нетипичного поведения заявок на закупку\\\large Краткий отчёт}
\author{}\date{}
\begin{document}
\maketitle\vspace{-2.2em}

\section{Инструменты и воспроизводимые шаги}
Python 3 (pandas, numpy, scikit-learn) в Google Colab; ноутбук \texttt{anomaly\_analysis.ipynb} запускается целиком (\emph{Run all}), на вход берёт \texttt{students.csv}, на выходе даёт \texttt{anomalies.csv}. Шаги: (1) загрузка, UUID как строки, \texttt{utf-8-sig}, время в tz-aware UTC+03:00, сортировка по \texttt{ProcessInstanceId, EventIndex}; (2) проверки качества (\texttt{assert}); (3) таблица заданий: календарная и рабочая длительность (Пн--Пт 09:00--18:00, функция покрыта юнит-тестами: пятница$\to$понедельник, выходные, раннее утро); (4) маршрут каждого процесса как строка кодов этапов и результатов; (5) нормы по типам заявок; (6) правила аномалий; (7) контроль Isolation Forest (seed 42); (8) генерация объяснений из данных и валидация CSV.

\section{Что считается обычным поведением}
Журнал: 12\,162 события, 1000 процессов (standard 635, urgent 168, contract 197), 4591 задание. Норма определена по частоте маршрутов внутри типа заявки и по схеме из условия. Основной маршрут: проверка $\to$ руководитель $\to$ бюджет [$\to$ юрист для contract] $\to$ исполнение (standard 496 из 635, urgent 133 из 168, contract 150 из 197). Допустимые варианты: один возврат на доработку (119 процессов), отказ руководителя (51), открытая заявка. Роли по этапам однозначны (проверка~--- specialist, бюджет~--- finance, юрист~--- legal, исполнение~--- procurement, доработка~--- requester). Ни одно задание не создано вне рабочего календаря; рабочие длительности этапов имеют медианы от 0,14 ч (проверка urgent) до 5,2 ч (исполнение contract), urgent быстрее standard примерно в 3 раза на каждом этапе, это считаем свойством типа, а не аномалией.

\section{Критерии, сравнения и причины выбора}
\textbf{Нарушение}~--- противоречие обязательной схеме (отсутствует или переставлен этап, роль не соответствует этапу, лишняя проверка без возврата). \textbf{Необычное, но допустимое}~--- только статистический хвост без противоречия схеме; в список он попадает, лишь если резко отделён от остального распределения. Сравнения всегда внутри типа заявки и этапа (медиана, $p_{98}$, $p_{99}$ по \emph{рабочей} длительности). Порог для хвоста подбирался по разрыву в данных, а не заранее: возвратов у обычных заявок 0 или 1, затем сразу 5--8; длительность проверки заявки до 2,5 ч у всех, далее лишь значения ровно 8 и 60 рабочих часов; закрытые задания короче минуты есть только в 8 процессах. Открытое задание считается аномальным, если нижний предел ожидания больше $3\times p_{99}$ закрытых заданий того же этапа и типа.

\section{Результаты}
Всего \textbf{88 строк по 80 процессам} (8\,\% журнала), 10 типов:\\[2pt]
\begin{tabular}{@{}lcp{9.8cm}@{}}\toprule
Тип & Проц. & Критерий\\\midrule
Skipped manager approval & 16 & бюджет создан без этапа руководителя\\
Execution before budget approval & 8 & исполнение раньше бюджета (3 contract~--- ещё и раньше юриста)\\
Duplicate specialist check & 8 & две подряд проверки accept без возврата\\
Unplanned audit checks & 8 & пара проверок auditor вокруг бюджета, причина не видна\\
Manager task performed by requester & 8 & роль requester на этапе руководителя\\
Excessive rework loops & 8 & 5--8 возвратов при норме 0--1\\
Instant task chain & 8 & задания закрыты за секунды\\
Stalled specialist check & 8 & проверка ровно 60 рабочих ч (норма $\le$2 ч)\\
Slow specialist check & 8 & проверка ровно 8 рабочих ч у standard ($p_{98}=1{,}16$ ч)\\
Stalled open manager approval & 8 & открыто с янв.--фев., 900--1030 рабочих ч при $p_{99}\approx5$ ч\\\bottomrule
\end{tabular}\\[3pt]
Классы почти не пересекаются: 8 процессов с пропуском руководителя совпали с 60-часовой проверкой (по две строки).

\textbf{Примеры.} \texttt{5bd32a9a}: standard, согласование руководителем создано 21.01.2026 10:09 и не закрыто к срезу, не менее 1033,85 рабочих ч при медиане 2,00 ч и $p_{99}=5{,}02$ ч. \texttt{5d7cf7bb}: standard, 8 возвратов на доработку (28 заданий), у остальных 992 заявок возвратов не более одного. \texttt{26b4c1f5}: standard, все 4 задания (проверка, руководитель, бюджет, исполнение) закрыты за 7--19 секунд, медиана согласования руководителем 2,00 ч.

\section{Ложные срабатывания, ограничения, ручная проверка}
\textbf{Данные.} Структурно журнал чистый (индексы, парность заданий, согласованность результатов проверены). Найден дефект: 25 событий в 10 процессах имеют отметки 1 июля 09:16--17:59 (UTC+03:00), то есть позже заявленного среза 00:00. Это завершённые процессы, не совпадающие с 10 открытыми. Мы считаем это дефектом выгрузки: события оставлены в данных с флагом, сами по себе аномалией не считаются. Для нижнего предела открытых заданий использован заявленный срез; расхождение не превышает суток и не влияет на выводы (ожидание порядка 1000 ч).

\textbf{Открытые заявки.} Всего 10. Восемь зависших согласований включены (ожидание превышает норму более чем в 100 раз); два исполнения от 30 июня (0,08 и 3,85 рабочих ч) в пределах нормы и не включены.

\textbf{Для ручной проверки.} (а) \emph{Unplanned audit checks}: внеплановая проверка допустима, а редкость сама по себе не доказывает нарушения; включены, так как 8 процессов выглядят одинаково (пара accept до и после бюджета), по данным причина не видна, при положительном ответе бизнеса строки можно исключить. (б) Исполнение standard 17,5 и 17,3 рабочих ч (максимум остальных 14,5 ч): плавный хвост без разрыва, не включено. (в) Единый инициатор у всех заявок: проверка <<согласовал сам себе>> по данным малоинформативна.

\textbf{Контроль.} Isolation Forest (500 деревьев, 14 признаков процесса, нормировка по типу) использован как независимая проверка. Он уверенно подтверждает возвраты (медианный ранг флагов 4 из 1000), открытые зависшие (14), 60-часовые проверки (26), мгновенные цепочки и внеплановые проверки (по 32), но почти не видит нарушения порядка и ролей (ранги 170--250) и 8-часовую проверку (521). Просмотр верхних процессов без правил показал лишь обычные короткие отказы и маршруты с одним возвратом, новых аномалий он не добавил. Нейросеть не применялась: нет меток, 1000 объектов, а требуются проверяемые объяснения.

\section{Использование AI}
Claude (Anthropic) использовался для разбора условия и плана анализа, написания кода ноутбука и текстов объяснений. Все числа в объяснениях вычисляются кодом из журнала, а не набраны вручную. Проверка: \texttt{assert}-проверки данных и итогового CSV (заголовки, существующие ID, отсутствие дублей), юнит-тесты функции рабочего времени, перекрёстная проверка Isolation Forest, полный перезапуск ноутбука с идентичным результатом. % TODO автору: допишите, что именно вы проверяли вручную (например, трассы выборочных процессов), либо удалите это замечание.
\end{document}
